%!TEX root = ../sztuthesis_main.tex

\section{实验设计与评估}

\subsection{问题定义}
本章目标是验证平台是否能够稳定支撑“参数配置-运行-采样-导出-分析”的实验流程，并观察无线参数变化对 BER、无线处理时延与轨迹误差的影响趋势。根据项目说明，当前数据属于仿真统计，因此本章侧重流程可复现性与趋势分析模板，而非宣称真实公网性能结论。

\subsection{实验目的}
（1）验证系统在 basic\_awgn 与 advanced\_cdl\_ofdm 模式下可稳定运行；

（2）验证 Eb/No 变化对 BER 与轨迹误差的影响方向；

（3）验证 CSV 导出字段能够支持离线统计和绘图。

\subsection{实验环境}
\begin{table}[htbp]
\centering
\caption{实验环境配置表}
\label{tab:exp-env}
\begin{tabular}{ll}
\hline
项目 & 配置 \\
\hline
操作系统 & Windows（开发机） \\
容器编排 & Docker Compose \\
后端服务 & FastAPI + Uvicorn \\
前端服务 & Vite 构建 + Nginx 静态托管 \\
物理引擎模式 & lightweight（默认）/ mujoco（可选） \\
\hline
\end{tabular}
\end{table}

\subsection{变量与指标}
\subsubsection{自变量}
（1）无线模式：basic\_awgn、advanced\_cdl\_ofdm；

（2）Eb/No（dB）：按实验组配置（待补充）；

（3）命令类型：robot\_joint\_control、robot\_ee\_control。

\subsubsection{观测指标}
（1）average\_ber、last\_ber；

（2）wireless\_delay\_ms、total\_wireless\_delay\_ms；

（3）trajectory\_error\_mean；

（4）queue\_pending；

（5）transmission\_count 与 total\_run\_time\_s。

\subsection{实验步骤}
（1）启动系统并确认 /health 正常；

（2）在前端设置无线模式与 Eb/No；

（3）发送连续控制命令并维持固定实验时长；

（4）导出 CSV 数据并进行离线统计；

（5）重复不同参数组并进行对比。

\begin{figure}[htbp]
\centering
\fbox{\parbox{0.92\textwidth}{
\textbf{图位占位：实验流程图。}\\
配置参数、发送命令、采样、导出、统计分析形成闭环。}}
\caption{实验执行流程示意}
\label{fig:exp-flow}
\end{figure}

\noindent 本图流程定义已整理至流程图页面（diagrams/flowcharts.html）中对应条目“fig:exp-flow”，可在浏览器查看并导出为 SVG/PNG 后替换本图位。
\noindent 图意说明：用于保证实验过程可复现。

\subsection{结果模板（待填）}
由于仓库中未提供固定批次实验原始数据文件，以下结果以模板方式给出，禁止填入未经验证数字。

\begin{table}[htbp]
\centering
\caption{不同无线模式与 Eb/No 的实验结果模板（待填）}
\label{tab:exp-result-template}
\begin{tabular}{llllll}
\hline
组别 & 模式 & Eb/No(dB) & average\_ber & wireless\_delay\_ms & trajectory\_error\_mean \\
\hline
G1 & basic\_awgn & 待补充 & 待补充 & 待补充 & 待补充 \\
G2 & basic\_awgn & 待补充 & 待补充 & 待补充 & 待补充 \\
G3 & advanced\_cdl\_ofdm & 待补充 & 待补充 & 待补充 & 待补充 \\
G4 & advanced\_cdl\_ofdm & 待补充 & 待补充 & 待补充 & 待补充 \\
\hline
\end{tabular}
\end{table}

\begin{table}[htbp]
\centering
\caption{CSV 字段核对表}
\label{tab:exp-csv-fields}
\begin{tabular}{ll}
\hline
字段名 & 用途 \\
\hline
average\_ber / last\_ber & 误码率统计 \\
wireless\_delay\_ms & 单次无线处理时延 \\
trajectory\_error\_mean & 轨迹偏差评估 \\
queue\_pending & 队列积压观测 \\
wireless\_mode / ebno\_db & 分组对比条件 \\
\hline
\end{tabular}
\end{table}

\begin{figure}[htbp]
\centering
\fbox{\parbox{0.92\textwidth}{
\textbf{图位占位：实验指标分析图。}\\
建议绘制 BER-Eb/No 曲线、时延箱线图、轨迹误差对比柱状图。}}
\caption{实验结果可视化建议（占位）}
\label{fig:exp-analysis-template}
\end{figure}

\noindent 本图流程定义已整理至流程图页面（diagrams/flowcharts.html）中对应条目“fig:exp-analysis-template”，可在浏览器查看并导出为 SVG/PNG 后替换本图位。
\noindent 插图位置建议：图~\ref{fig:exp-flow} 放在步骤后，图~\ref{fig:exp-analysis-template} 放在结果模板后。

\subsection{结果分析与局限性}
预期可验证结论模板如下：

（1）在其他条件近似相同情况下，Eb/No 提升通常对应 BER 下降趋势；

（2）advanced\_cdl\_ofdm 模式在 BER 上可能更优，但处理时延可能更高；

（3）队列积压上升会增加控制链路整体滞后感。

局限性说明：

（1）当前实验基于仿真模型，不等价于真实公网链路；

（2）未包含大规模并发用户场景压测数据（待补充）；

（3）未提供跨地区网络实测样本（待补充）。

\subsection{本章小结}
本章给出了结构化实验框架、变量指标定义、执行步骤和结果模板，确保后续补充数据时可直接落地，不需要重写论文结构。

\subsection{事实核对清单}
（1）已核对：指标命名与后端 telemetry/csv 字段一致；

（2）已核对：项目支持模式切换与 Eb/No 参数化；

（3）已核对：CSV 导出接口可用于离线分析；

（4）待补充：真实实验数值与统计显著性分析结果。
